%% kapa: exam=1 round=26 year=2015 period=1 form=A subject=civil_law subject_ko=민법 seq=1 number=1
%% source: Q-Net 감정평가사 기출문제 (cst00309) artlSeq=5138001 file=extracted/1008/5138001/제26회_감정평가사 1차_기출/제26회_감정평가사 1차_기출/1교시_감평_A형.pdf
%% conversion: kordoc 4.24.51 -> tools/gen1.py (auto); review=proofread-sonnet; flags=none
\begin{question}{1}
관습법과 사실인 관습에 관한 설명으로 옳지 않은 것은? (다툼이 있으면 판례에 따름)
\begin{choices}
  \item 관습법이 되기 위해서는 사회구성원의 법적 확신이 필요하다.
  \item 사실인 관습은 법령으로서의 효력이 없는 단순한 관행에 지나지 않으므로 법률행위 당사자의 의사를 보충하는 수단이 될 수 없다.
  \item 여성은 종중의 구성원이 될 자격이 없다는 종래의 관습에 법적 효력은 인정되지 않는다.
  \item 온천에 관한 권리는 관습법상 권리가 아니다.
  \item 관습법의 존부는, 법원이 알 수 없는 경우를 제외하고는, 당사자의 주장ㆍ증명을 기다리지 않고 법원이 직권으로 이를 확정하여야 한다.
\end{choices}
\end{question}
\answer{2}
